Fix \(p\in(0,1)\), \(\alpha\in(0,1)\), and one deterministic allocation sequence \((n_s(n),n_t(n))_{n\ge2}\) satisfying Assumption~\ref{ass:triangular-allocation}.  Thus \(n_s(n),n_t(n)\ge1\), \(n_s(n)+n_t(n)=n\) for every \(n\ge2\), and \(n_s(n)/n\to p\).  Every model and array below is sampled using this same sequence.  For each row, the sample space is the common finite product
\[
\Omega_n=\Omega(n_s(n),n_t(n))
=\bigl(\mathcal E\times\mathcal W\times\mathcal X\times\mathcal Y\bigr)^{n_s(n)}
\times\mathcal W^{n_t(n)},
\tag{1}
\]
equipped with its power-set sigma-field.  The updated row-wise sampling assumptions identify the row law on (1) with the finite product measure \(P_{O,n}^{\otimes n_s(n)}\otimes b_n^{\otimes n_t(n)}\); no infinite i.i.d. sequence is used.
Every row belongs to \(\mathfrak M_{+}\), so strict primitive positivity makes every one-observation source and target cell positive.  Thus (1) is also the common full support of every row law in the class.

\paragraph{Fixed regular baseline.}
Let
\[
\mathcal E=\{-1,+1\},\qquad
\mathcal U=\mathcal W=\mathcal Y=\{0,1\},\qquad
\mathcal X=\{x\}.
\]
Order binary vectors as \((1,0)\), environment columns as \((+1,-1)\), and latent columns of the proxy channel as \((1,0)\).  Define \(\mathcal M^{(0)}\) by
\[
\pi=(1/2,1/2),\qquad
S^{(0)}=\begin{pmatrix}1/2&1/2\\1/2&1/2\end{pmatrix},\qquad
M^{(0)}=\begin{pmatrix}3/4&1/4\\1/4&3/4\end{pmatrix},
\]
\[
a_x(u)=1,\qquad q^{(0)}=(1/2,1/2)^{\top},\qquad
f^{(0)}_{x,1}(u,w)=f^{(0)}_{x,0}(u,w)=1/2.
\tag{2}
\]
Every primitive entry is positive and stochastic, and \(\det(M^{(0)})=1/2\ne0\).  The factorization, invariance, positivity, and proxy-injectivity assumptions imply \(\mathcal M^{(0)}\in\mathfrak M_{+}\).

Summing the source factorization over \(u\) gives
\[
P_{O,0}(e,w,x,y)=\frac18
\tag{3}
\]
for every \(e,w,y\).  Target invariance gives \(b_0=(1/2,1/2)^{\top}\) and \(\theta_{x,1}(\mathcal M^{(0)})=1/2\).  Therefore
\[
H_{x,0}=\begin{pmatrix}1/4&1/4\\1/4&1/4\end{pmatrix},
\qquad z_{x,1,0}=(1/4,1/4)^{\top}.
\tag{4}
\]
The singular values of \(H_{x,0}\) are \(1/2\) and \(0\), so
\[
\operatorname{rank}(H_{x,0})=1,
\quad \sigma_1(H_{x,0})=1/2\ge1/4,
\quad b_0=H_{x,0}(1,1)^{\top}.
\tag{5}
\]
All source observed cells are \(1/8\), and all target proxy cells are \(1/2\), so the cell floor is at least \(1/16\).

For the variance condition, rank one in (4) gives
\[
H_{x,0}^{\dagger}=\begin{pmatrix}1&1\\1&1\end{pmatrix},
\qquad
D_z\phi_1(H_{x,0},z_{x,1,0},b_0)
=\bigl(H_{x,0}^{\dagger}b_0\bigr)^{\top}=(1,1).
\tag{6}
\]
The source part of the scalar influence variable is therefore \(\mathbf 1\{Y=1\}+\psi(E,W)\), up to an additive constant, for a deterministic \(\psi\) contributed by the \(H\)-coordinates.  Under (2), \(Y\) is conditionally fair given \((E,W)\).  For any finite random variable \(T\), expanding \(\mathbb E[T^2]-\mathbb E[T]^2\) and adding and subtracting \(\mathbb E[\mathbb E(T\mid E,W)^2]\) proves the conditional variance identity.  Applying that algebra here yields
\[
\begin{aligned}
\operatorname{Var}_{P_{O,0}}\!\left(\mathbf 1\{Y=1\}+\psi(E,W)\right)
&=\mathbb E\!\left[\operatorname{Var}(\mathbf 1\{Y=1\}\mid E,W)\right]\\
&\quad+\operatorname{Var}\!\left(\mathbb E[\mathbf 1\{Y=1\}+\psi(E,W)\mid E,W]\right)\\
&\ge\frac14.
\end{aligned}
\tag{7}
\]
The source allocation factor is \(p^{-1}\), and the independent target contribution is nonnegative.  Hence
\[
V_{1,x,1}(\mathcal M^{(0)};p)\ge\frac1{4p}>\frac1{16}.
\tag{8}
\]
Equations (3)--(8) verify all five conditions of Definition~\ref{def:strongly-identified-submodel}; thus \(\mathcal M^{(0)}\in\mathcal R_{1,1/4,1/16}\).

\paragraph{Admissible weak-rank array.}
For \(n\ge2\), put \(\gamma_n=\delta_n=n^{-2}\) and define \(\mathcal M^{(1)}_n\) with the same \(\pi,a_x\) and
\[
S^{(1)}_n=\begin{pmatrix}1/2+\gamma_n&1/2-\gamma_n\\1/2-\gamma_n&1/2+\gamma_n\end{pmatrix},
\qquad
M^{(1)}_n=\begin{pmatrix}1/2+\delta_n&1/2-\delta_n\\1/2-\delta_n&1/2+\delta_n\end{pmatrix},
\]
\[
q^{(1)}=(3/4,1/4)^{\top},\qquad
f^{(1)}_{x,1}(1,w)=3/4,\qquad
f^{(1)}_{x,1}(0,w)=1/4,\qquad
f^{(1)}_{x,0}=1-f^{(1)}_{x,1}.
\tag{9}
\]
For \(n\ge2\), all entries are positive, and
\[
\det M^{(1)}_n=2\delta_n>0,
\qquad
\det S^{(1)}_n=2\gamma_n>0.
\tag{10}
\]
Thus \(\mathcal M^{(1)}_n\in\mathfrak M_{+}\) in every row.

Because \(x\) is the sole treatment value,
\[
B^{(1)}_{x,n}=M^{(1)}_nS^{(1)}_n
=\begin{pmatrix}
1/2+2\delta_n\gamma_n&1/2-2\delta_n\gamma_n\\
1/2-2\delta_n\gamma_n&1/2+2\delta_n\gamma_n
\end{pmatrix},
\tag{11}
\]
\[
H^{(1)}_{x,n}=\frac12B^{(1)}_{x,n},
\qquad
\det B^{(1)}_{x,n}=4\delta_n\gamma_n>0.
\tag{12}
\]
Target invariance gives
\[
b^{(1)}_n=M^{(1)}_nq^{(1)}
=\begin{pmatrix}1/2+\delta_n/2\\1/2-\delta_n/2\end{pmatrix}.
\tag{13}
\]
By (12), \(b^{(1)}_n\in\operatorname{col}(B^{(1)}_{x,n})\) for every \(n\).  The array is assigned the one fixed allocation sequence chosen before (1), and the two finite coordinate blocks have the stipulated product law on \(\Omega_n\) by the row-wise sampling assumptions.  Hence Definition~\ref{def:studentized-array-class} gives
\[
\mathbb M^{(1)}=(\mathcal M^{(1)}_n)_{n\ge2}\in\mathcal A_p.
\tag{14}
\]
If \(\lambda_n\) is the unique solution of \(B^{(1)}_{x,n}\lambda_n=b^{(1)}_n\), addition and subtraction of its two row equations give
\[
\lambda_{1,n}+\lambda_{2,n}=1,
\qquad
\lambda_{1,n}-\lambda_{2,n}=\frac1{4\gamma_n}.
\tag{15}
\]
Thus the balancing norm diverges, including after conversion to unconditional weights because \(P_{O,n}(E=e,X=x)=1/2\).

The causal value is separated from the baseline:
\[
\theta_{x,1,n}(\mathcal M^{(1)}_n)
=\sum_{u,w}q^{(1)}_uM^{(1)}_n(w,u)f^{(1)}_{x,1}(u,w)
=\frac34\frac34+\frac14\frac14
=\frac58.
\tag{16}
\]

\paragraph{Indistinguishability on the common row spaces.}
For finite probability vectors on a common space, write \(\lVert P-Q\rVert_{\mathrm{TV}}=\frac12\sum_a|P(a)-Q(a)|\).  Retain \(f^{(1)}\), but replace \(S^{(1)}_n,M^{(1)}_n\) by the all-\(1/2\) binary kernels \(S^{\star},M^{\star}\).  The resulting observed law equals \(P_{O,0}\): \(E,W\) are independent and fair, and averaging response probabilities \(3/4,1/4\) over fair \(U\) makes \(Y\) fair independently of \((E,W)\).

At a full-data cell, adding and subtracting the density with \(S^{\star}\) and \(M^{(1)}_n\) gives
\[
\begin{aligned}
&\pi_eS^{(1)}_n(u,e)M^{(1)}_n(w,u)f^{(1)}_{x,y}(u,w)
-\pi_eS^{\star}(u,e)M^{\star}(w,u)f^{(1)}_{x,y}(u,w)\\
&=\pi_e\bigl(S^{(1)}_n(u,e)-S^{\star}(u,e)\bigr)
M^{(1)}_n(w,u)f^{(1)}_{x,y}(u,w)\\
&\quad+\pi_eS^{\star}(u,e)
\bigl(M^{(1)}_n(w,u)-M^{\star}(w,u)\bigr)
f^{(1)}_{x,y}(u,w).
\end{aligned}
\]
After summing absolute values over all cells, normalization of the unchanged kernels reduces the two terms to the conditional binary-kernel total variations \(\gamma_n\) and \(\delta_n\).  Marginalization groups signed summands before taking absolute values and hence cannot increase total variation.  Therefore
\[
\lVert P^{(1)}_{O,n}-P_{O,0}\rVert_{\mathrm{TV}}
\le\gamma_n+\delta_n=2n^{-2},
\qquad
\lVert b^{(1)}_n-b_0\rVert_{\mathrm{TV}}=\delta_n/2.
\tag{17}
\]
For completeness, the product-measure inequality follows from the telescoping identity
\[
\bigotimes_{i=1}^N P_i-\bigotimes_{i=1}^N Q_i
=\sum_{j=1}^N
\left(\bigotimes_{i<j}P_i\right)\otimes(P_j-Q_j)\otimes
\left(\bigotimes_{i>j}Q_i\right)
\tag{18}
\]
and the triangle inequality, because probability vectors have \(\ell_1\)-norm one.  Both product laws below are measures on the same \(\Omega_n\) from (1), and (17)--(18) give
\[
\begin{aligned}
&\left\lVert
(P^{(1)}_{O,n})^{\otimes n_s(n)}\otimes(b^{(1)}_n)^{\otimes n_t(n)}
-P_{O,0}^{\otimes n_s(n)}\otimes b_0^{\otimes n_t(n)}
\right\rVert_{\mathrm{TV}}\\
&\qquad\le n_s(n)(\gamma_n+\delta_n)+n_t(n)\delta_n/2
\le2n^{-1}\longrightarrow0.
\end{aligned}
\tag{19}
\]

\paragraph{Transferred coverage.}
Let \(\mathcal C_n:\Omega_n\to\{C\subseteq[0,1]:C\ne\varnothing\}\) be any sequence satisfying the theorem's uniform-coverage premise, and put \(A_n=\{5/8\in\mathcal C_n\}\).  Since \(\Omega_n\) is finite with its power-set sigma-field, \(A_n\) is measurable.  Because every law in the infimum over \(\mathcal A_p\) is a law on this same \(\Omega_n\), the premise and (14) imply
\[
\liminf_{n\to\infty}
\mathbb P_{\mathcal M^{(1)}_n}^{\,n_s(n),n_t(n)}(A_n)
\ge1-\alpha.
\tag{20}
\]
For any event, the difference of its probabilities under two finite laws is bounded by their total variation distance.  Equations (19)--(20) therefore give
\[
\liminf_{n\to\infty}
\mathbb P_{\mathcal M^{(0)}}^{\,n_s(n),n_t(n)}(A_n)
\ge1-\alpha.
\tag{21}
\]

\paragraph{Localization of the fixed-model Wald interval.}
Under \(\mathcal M^{(0)}\), each coordinate of \(\widehat H_{x,n},\widehat z_{x,1,n},\widehat b_n\) is an average of independent indicators.  If \(\widehat\mu\) averages \(N\) indicators, then
\[
\mathbb E[(\widehat\mu-\mu)^2]\le\frac1{4N},
\qquad
\Pr\{|\widehat\mu-\mu|>t\}\le\frac1{4Nt^2}.
\tag{22}
\]
There are finitely many coordinates; the fixed allocation assumption gives \(n_s(n)/n\to p\) and \(n_t(n)/n\to1-p\), both positive.  A finite union bound yields
\[
(\widehat H_{x,n},\widehat z_{x,1,n},\widehat b_n)
-(H_{x,0},z_{x,1,0},b_0)=O_P(n^{-1/2}).
\tag{23}
\]

At \(H_{x,0}\), the eigenvalues of \(H_{x,0}H_{x,0}^{\top}\) are \(1/4\) and \(0\).  In a neighborhood of \(H_{x,0}\), the roots of the quadratic characteristic polynomial remain separated, and the leading projector is
\[
P_1(H)=\frac{HH^{\top}-\lambda_2(H)I}{\lambda_1(H)-\lambda_2(H)}.
\tag{24}
\]
Thus the roots, \(P_1(H)\), the rank-one truncation, its pseudoinverse, \(\phi_1\), and \(D\phi_1\) are continuously differentiable with bounded derivatives on a smaller neighborhood.  The plug-in multinomial covariance is polynomial in empirical cell frequencies and bounded on the finite simplex.  The mean-value inequality, (23), and the definitions of the Wald estimator and totalized variance give
\[
\widehat\theta^{\,W}_{x,1,n}=\frac12+O_P(n^{-1/2}),
\qquad
\widehat V_n=O_P(1).
\tag{25}
\]
Since the normal quantile is finite, the untruncated Wald radius is \(O_P(n^{-1/2})\), and the clipped interval is nonempty with probability tending to one.  Put
\[
\Delta_n=
\begin{cases}
\sup_{\zeta\in I^W_{n,1-\alpha}}|\zeta-1/2|,
&I^W_{n,1-\alpha}\ne\varnothing,\\
1,&I^W_{n,1-\alpha}=\varnothing.
\end{cases}
\]
Clipping cannot increase distance from \(1/2\), so
\[
\Delta_n=O_P(n^{-1/2})=o_P(1).
\tag{26}
\]

\paragraph{Quantitative impossibility.}
Fix \(\varepsilon\in(0,1/8)\) and let \(G_n=\{\Delta_n\le\varepsilon\}\).  Equation (26) gives \(\mathbb P_{\mathcal M^{(0)}}(G_n^c)\to0\).  On \(A_n\cap G_n\), the Wald interval is nonempty and \(5/8\in\mathcal C_n\), whence the first directed Hausdorff term gives
\[
\begin{aligned}
d_H(\mathcal C_n,I^W_{n,1-\alpha})
&\ge\inf_{\zeta\in I^W_{n,1-\alpha}}|5/8-\zeta|\\
&\ge\frac18-\sup_{\zeta\in I^W_{n,1-\alpha}}|\zeta-1/2|\\
&\ge\frac18-\varepsilon.
\end{aligned}
\tag{27}
\]
Therefore (21), (26), and the union bound imply
\[
\liminf_{n\to\infty}
\mathbb P_{\mathcal M^{(0)}}^{\,n_s(n),n_t(n)}
\left\{I^W_{n,1-\alpha}\ne\varnothing,\ 
d_H(\mathcal C_n,I^W_{n,1-\alpha})\ge\frac18-\varepsilon\right\}
\ge1-\alpha.
\tag{28}
\]
Because the Wald interval is empty with probability tending to zero, the same bound holds after any totalization on empty samples.  Since \(1/8-\varepsilon>0\), (28) rules out unscaled Hausdorff localization and therefore also \(\sqrt n\)-scaled localization.

Finally, uniform coverage alone remains attainable.  For every \(\mathbb M\in\mathcal A_p\) and every \(n\ge2\), the span condition gives \(b_n\in\operatorname{col}(B_{x,n})\), while the fixed allocation assumption gives positive finite \(n_s(n),n_t(n)\).  Applying Theorem~\ref{thm:finite-sample-projection-coverage} to the row-n product law on \(\Omega_n\) gives coverage at least \(1-\alpha\) in every row.  Taking the infimum over the common-space class \(\mathcal A_p\), then the lower limit, preserves this bound even along changing-rank arrays.